\chapter{基于模型的知识获取}
\label{chap:knowledge-acquisition-model-based}

人工可以直接给出答案，规则可以把明确条件转成答案；但当待处理样本数量很大、输入形态复杂，或者答案难以写成有限规则时，这两条路径都可能遇到规模瓶颈。已经训练好的模型提供了第三种来源：它把过去数据中学到的关系、调用时给出的参考样本以及样本之间的结构，转化为当前任务的类别、数值、对象关系或文本答案。本章把这些输出统称为\term{候选知识}，强调它们可以进入后续训练或分析，但尚未因为出自模型而成为可靠事实。

要把模型输出真正变成可管理的知识来源，不能只执行一次推断。模型是否适合任务，取决于它接受什么输入、产生什么输出以及在哪些数据范围内经过检查；同一个模型看到不同任务说明、示范或外部证据，也可能给出不同答案；输出还可能通过近邻检索或图上传播获得，而不是只来自模型参数。反馈到来以后，需要进一步判断应修改模型、输入内容还是获取安排。由此，本章围绕一个问题展开：怎样准备并调用模型获得任务所需的候选知识，又怎样在不混淆证据用途的前提下改进整个获取过程？

花卉识别将作为贯穿案例。目标不是训练一个尽可能大的视觉模型，而是为一批照片补充物种名称、必要的弃权标记和可追溯的判断依据。模型可能是从少量标注照片训练的分类器，也可能是已有的花卉识别模型或接受图文指令的通用模型。无论来源如何，每条输出都要能够追溯到模型版本、输入版本和预测机制，后续复核才知道错误从哪里进入。

本章先说明模型怎样得到花卉答案，再讨论这些答案怎样参与下一轮训练。记录沿用定义~\ref{def:knowledge-acquisition-record}的约定，各节只补充当前机制特有的输入与依赖。

\section{标注模型的准备}
\label{sec:knowledge-acquisition-model-preparation}

\term{标注模型}（labeling model）是本章对“承担候选知识获取任务的模型”的功能性称呼，并不限定模型结构。准备这种模型时，需要依次回答三个问题：有哪些可试验的方案，哪些方案需要训练或适配，以及小规模试标能否支持扩大使用范围。方案选择给出候选模型与资源约束，训练适配形成可调用版本，试标确认则给出已经检查过的能力范围和后续行动。

模型来源与准备步骤不是一一对应的。已有专用模型可能因类别体系变化而需要适配，通用基础模型可能在任务说明和示范足够清楚时无需更新参数，从头构建的模型也可能因标签不足而无法进入训练。图~\ref{fig:knowledge-acquisition-model-preparation}把这些共同步骤放在一条流程中；反馈箭头表明，试标失败以后应回到具体成因，而不是默认把模型换得更大。

\begin{figure}[htbp]
  \centering
  % 图4-1：标注模型的准备与试标决策，按112 mm正文宽度设计。
\begin{tikzpicture}[
  x=1mm,
  y=1mm,
  font=\figurefont\footnotesize,
  text=BookInk,
  prep input/.style={book node,draw=FigureInput,fill=FigureInput!7,line width=0.8pt,text width=25mm,minimum height=11mm,inner sep=1.5mm},
  prep process/.style={book node,draw=FigureModel,fill=FigureModel!7,line width=0.8pt,text width=28mm,minimum height=11mm,inner sep=1.5mm},
  prep decision/.style={book node,draw=BookAmber,fill=BookAmber!7,line width=0.9pt,text width=24mm,minimum height=11mm,inner sep=1.5mm},
  prep result/.style={book node,draw=BookMuted,fill=BookPaper,line width=0.7pt,text width=25mm,minimum height=11mm,inner sep=1.5mm},
  prep flow/.style={-{Stealth[length=2mm]},draw=BookMuted,line width=0.8pt,rounded corners=1mm},
  prep feedback/.style={-{Stealth[length=2mm]},draw=BookAmber,line width=0.8pt,dashed,rounded corners=1mm},
  edge text/.style={fill=white,inner sep=0.8pt,font=\figurefont\scriptsize,text=BookMuted}
]
  \node[prep input] (requirements) at (16,0) {任务要求\\类别、边界、弃权};
  \node[prep input] (models) at (55,0) {可用模型\\专用或通用};
  \node[prep input] (labels) at (94,0) {初始标签\\数据与覆盖};

  \node[prep process] (select) at (55,-19) {选择候选方案\\并记录资源约束};
  \node[prep decision] (need) at (55,-37) {是否需要\\参数训练？};
  \node[prep process] (adapt) at (87,-51) {训练或适配\\形成模型版本};
  \node[prep process] (pilot) at (55,-66) {试标确认\\能力范围与失败类型};

  \node[prep result,draw=BookAmber] (source) at (17,-86) {补充其他来源\\或重定任务};
  \node[prep result,draw=FigureOutput] (accept) at (55,-86) {采用当前方案\\进入批量获取};
  \node[prep result,draw=FigureModel] (continue) at (93,-86) {继续适配\\再做试标};

  \draw[prep flow] (requirements.south) -- (16,-10) -| (select.west);
  \draw[prep flow] (models.south) -- (select.north);
  \draw[prep flow] (labels.south) -- (94,-10) -| (select.east);
  \draw[prep flow] (select.south) -- (need.north);
  \draw[prep flow] (need.south) -- node[edge text,right] {否} (pilot.north);
  \draw[prep flow] (need.east) -- node[edge text,above] {是} (72,-37) |- (adapt.north);
  \draw[prep flow] (adapt.south) |- (pilot.east);
  \draw[prep flow] (pilot.south) -- (accept.north);
  \draw[prep flow] (pilot.west) -- (36,-66) |- (source.north);
  \draw[prep flow] (pilot.east) -- (75,-66) |- (continue.north);

  \draw[prep feedback] (source.west) -- (1,-86) |- (select.west);
  \draw[prep feedback] (continue.east) -- (110,-86) |- (adapt.east);
  \node[edge text,anchor=west] at (3,-29) {资源或要求变化};
  \node[edge text,anchor=east] at (108,-42) {新训练版本};
\end{tikzpicture}

  \caption{标注模型的准备与试标决策。任务要求、可用模型和初始标签共同约束方案选择；不需要参数训练的方案可以直接进入试标。试标只支持准备阶段的采用、继续适配或补充来源决策，不替代方案冻结后的独立验收。}
  \label{fig:knowledge-acquisition-model-preparation}
\end{figure}

\subsection{模型方案的选择}
\label{sec:knowledge-acquisition-model-scheme-selection}

模型方案应从任务接口和可用资源出发，而不是从模型名称出发。花卉识别至少要明确输入是否只有照片，还是还包括拍摄地点、季节和植株部位；输出是封闭物种表中的单个类别、多个候选，还是允许返回“未知”并记录“证据不足”原因；覆盖范围是本地常见花卉，还是开放世界中的任意植物。只有这些条件固定以后，才能判断现有模型的输入模态、类别体系和输出结构是否相容。

从准备动作看，可以比较三种方案。若只有任务数据和标签，就要从头构建模型；若已有模型的任务语义和输入条件与当前需求一致，可以直接复用；若已有参数提供了有用表示，但输出类别、数据分布或接口不完全一致，则需要适配。专用模型和通用模型都可能落入后两种方案，规模也不能代替适用性判断。模型卡一类报告框架之所以要求记录预期用途、评价条件和已知限制，正是因为“已有模型”本身不足以说明它能否在当前场景中使用\scite{mitchell2019modelcards}

表~\ref{tab:knowledge-acquisition-model-preparation}把三个准备方案放在同一资源口径下。一次性投入包括标签建设、训练、接口改造和试验；持续条件则包括本地推断资源、外部调用限制、索引或版本维护。两类成本不能相互替代：直接调用模型可能减少初始训练，却把数据传输、限额和版本变化带入每一轮获取；从头构建的模型有较高前期投入，却可能允许更细地控制输入、参数和部署边界。直接复用时通常只能控制调用配置和输出约束，适配时还要明确哪些参数可以改变。

\begin{table*}[htbp]
  \centering
  \small
  \setlength{\RaggedRightParindent}{0pt}
  \renewcommand{\arraystretch}{1.15}
  \begin{tabularx}{\linewidth}{@{}>{\RaggedRight\arraybackslash}p{20mm}*{4}{>{\RaggedRight\arraybackslash}X}@{}}
    \toprule
    准备方案 & 起点与准备 & 一次性投入 & 持续条件 & 重点试验 \\
    \midrule
    从头构建 & 任务标签；训练并接入模型 & 标签、训练与调参 & 算力、维护与重训 & 标签覆盖、过拟合 \\
    直接复用 & 已有模型；对齐接口 & 接入、试标与合规检查 & 服务、限额与版本 & 语义错配、接口漂移 \\
    适配已有模型 & 预训练参数；确定更新范围 & 适配训练与能力检查 & 推断资源与再适配 & 负迁移、原能力下降 \\
    \bottomrule
  \end{tabularx}
  \caption{模型准备方案的资源与适用条件。三行描述准备动作，不是互斥的模型类别；同一个已有模型可以直接复用，也可以在类别或输入条件变化后转入适配。成本列给出需要测量的构成，不预设具体价格或效果排序。}
  \label{tab:knowledge-acquisition-model-preparation}
\end{table*}

以花卉识别为例，少量照片已有物种标签时，可以训练轻量分类器或适配视觉编码器；存在经过同一物种体系训练的识别模型时，可以先直接试标；只有通用多模态模型时，则要通过任务说明、示范和输出约束检查它能否稳定使用本项目的类别表。三种资源条件都不能跳过困难物种、照片质量和类别体系外样本的试验，尤其要检查少数类别缺失、输入变化和高置信错误。选择阶段得到的是带有资源约束和假定适用范围的候选方案，是否采用仍由试标决定。

\subsection{标注模型的训练与适配}
\label{sec:knowledge-acquisition-model-training-adaptation}

只有在当前方案需要改变参数时，训练才是准备步骤。设带任务答案的训练集为$\mathcal{D}_{\mathrm{tr}}=\{(\bm{x}_i,y_i)\}_{i=1}^{n}$，模型参数为$\bm{\theta}$，需要保留的预训练参数为$\bm{\theta}_0$。一种概括性的适配目标是
\begin{equation}
  \widehat{\bm{\theta}}
  \in
  \argmin_{\bm{\theta}\in\Theta_{\mathrm{upd}}}
  \frac{1}{n}\sum_{i=1}^{n}\ell\!\left(f_{\bm{\theta}}(\bm{x}_i),y_i\right)
  +\lambda\Omega(\bm{\theta},\bm{\theta}_0),
  \label{eq:knowledge-acquisition-adaptation-objective}
\end{equation}
其中$\Theta_{\mathrm{upd}}$规定允许更新的参数，$\ell$衡量任务输出与参照答案的差异，$\Omega$用于表达保持已有参数或能力的偏好，$\lambda\geq0$控制其强度。从头训练时可以令$\Theta_{\mathrm{upd}}$包含全部参数并去掉相对于$\bm{\theta}_0$的约束；只训练输出层时，表示部分被冻结；微调则允许部分或全部预训练参数改变。式~\eqref{eq:knowledge-acquisition-adaptation-objective}只概括信息怎样进入训练，不声称任意正则项都能防止原能力下降。

训练数据要按获取对象而不是只按文件随机划分。若同一植株的连拍照片分别进入训练集、方案选择集和最终检查集，模型可能借助几乎相同的背景和外观通过检查，却没有学会识别新的植株。更稳妥的做法是按植株或采集批次分组划分，再检查各物种、拍摄条件和语义状态为未知的情形是否得到覆盖。初始标签错误会直接改变损失目标，少数类别缺失则使模型没有机会学到相应输出；这些问题不能靠延长训练自动消失。

训练过程还要明确停止和回退依据。训练损失用于观察优化是否按目标进行，方案选择集用于比较冻结范围、权重和停止时刻，最终检查集则在这些决定冻结后使用。若适配后常见花卉识别有所改善，却丢失了原模型对低清照片的弃权能力，就不能只报告平均正确率。模型版本应与训练数据版本、标签体系、参数更新范围和推断配置一同保存，使后续候选知识能够追溯到实际调用的版本。

已有模型已经满足任务接口和试标要求时，可以跳过参数训练。此时仍要冻结具体模型版本和调用配置，但不能为了让流程看起来完整而虚构一次训练。相反，若输出类别或输入模态根本不相容，少量微调也未必足以补足能力；应回到表~\ref{tab:knowledge-acquisition-model-preparation}重新选择方案，而不是把所有失败都解释成优化不足。

\subsection{标注能力的试验与确认}
\label{sec:knowledge-acquisition-model-pilot-validation}

批量获取之前，需要用小规模\term{模型试标}（model pilot annotation）检查模型能否理解任务、返回可解析结果，并在已声明的范围内给出可用候选。它与第~\ref{sec:knowledge-acquisition-human-pilot}节的人工试标都属于小规模试运行，但执行主体、需要保存的记录和失败去向不同。模型试标样本应覆盖典型情形、易混边界、低质量输入、类别体系外对象和允许弃权的情形。对每条记录，至少保存样本标识、模型与输入版本、原始输出、解析状态、候选答案以及参照答案；调用失败、格式无效和语义错误应分别计数，因为它们需要不同修正。

试标的结论应落实为三种行动。若模型在预定范围内满足进入批量阶段的门槛，可以采用当前方案；若错误与可补充标签、可调整参数或接口有关，就继续适配；若关键证据不在模型可见输入中，或任务超出可确认范围，就补充人工、规则、外部文档等其他来源。试标集一旦反复用于选择模型、阈值和提示，就已经成为开发证据，不能再作为冻结方案的最终独立证明。具体的抽样、质量指标和区间估计由第~\ref{chap:knowledge-quality-assessment-correction}章处理。

% 长作者引文从段中开始，需同时预留段首至标记和完整边注的高度。
\Needspace{100mm}
当一个模型的输出用于监督另一个模型时，前者承担\term{任务教师}角色。这个名称只说明它提供任务目标，不保证教师更大、更准确，也不保证学生能力受教师上界严格限制。弱到强泛化研究在特定语言、棋类和奖励建模实验中观察到，强预训练模型用弱模型标签微调后可以超过弱监督者，但朴素微调仍未恢复强模型的全部能力\scite{burns2023weaktostrong}
因此，教师与学生的能力排序必须在当前任务和数据范围内检查，不能把该实验现象扩展成“学生总能纠正教师”的一般结论。

模型试标可以同时放入山茶与茶梅等相似物种、只有花瓣局部的照片，以及不在类别表中的植物。模型若只在典型正面照片上表现良好，适用范围就应收窄；若模型把应记为未知的样本强行映射到已有类别，需增加弃权接口并记录“信息不足”等原因，或补充其他来源。最终确认记录应同时写出可用范围、已知失败类型和人工接管条件，而不只保留一个总体分数。

\section{标注输入的组织与构造}
\label{sec:knowledge-acquisition-input-organization}

模型准备回答了“由谁作答”，下一步要回答“作答时看到了什么”。一次模型标注的输入通常承担四种功能：任务要求规定问题和输出含义，待处理样本及上下文提供当前观测，示范案例展示规范怎样落到实例上，辅助证据则补充外部事实或对象关系。四类内容可以组合，但不是每个模型都必须同时接收全部内容；固定接口模型也可以通过任务配置、标签映射和检索字段实现相同功能。

图~\ref{fig:knowledge-acquisition-input-composition}把“说明怎样作答”与“支持当前判断”区分开。任务要求和示范主要约束作答规则，当前样本和辅助证据主要支持本次判断；一条带答案参考照片可能同时承担示范和证据作用，此时必须记录它在当前调用中的具体用途，否则后续很难判断收益来自格式引导还是信息检索。

下面沿用同一个S042请求，在四个小节中依次加入任务、观测、示范和证据；将这些片段按顺序合并，就是一次可交给模型执行的完整请求。本例将照片的可见内容写成观测文本，以便仅凭书中文字检查输入是否充分。对象、名录、日期和输出均为教学构造，没有调用真实模型，也不作为植物鉴定依据。

\begin{figure}[htbp]
  \centering
  % 图4-2：一次模型标注中的四类输入，按112 mm正文宽度设计。
\begin{tikzpicture}[
  x=1mm,
  y=1mm,
  font=\figurefont\footnotesize,
  text=BookInk,
  input block/.style={book node,draw=FigureInput,fill=FigureInput!6,line width=0.8pt,text width=37mm,minimum height=12mm,inner sep=1.4mm},
  model block/.style={book node,draw=FigureModel,fill=FigureModel!8,line width=0.9pt,text width=25mm,minimum height=18mm,inner sep=1.5mm},
  output block/.style={book node,draw=FigureOutput,fill=FigureOutput!6,line width=0.8pt,text width=17mm,minimum height=18mm,inner sep=1.2mm},
  convention/.style={-{Stealth[length=2mm]},draw=BookAmber,line width=0.8pt,dashed},
  evidence/.style={-{Stealth[length=2mm]},draw=BookBlue,line width=0.9pt},
  edge text/.style={fill=white,inner sep=0.7pt,font=\figurefont\scriptsize,text=BookMuted}
]
  \node[input block,draw=BookAmber] (requirement) at (21,0) {任务要求与输出约定\\物种表、弃权、返回字段};
  \node[input block,draw=BookAmber] (demonstration) at (21,-18) {示范案例与参照答案\\S011：山茶；S017：弃权};
  \node[input block] (sample) at (21,-36) {待处理样本与必要上下文\\S042：花、叶、地点、时间};
  \node[input block] (support) at (21,-54) {辅助证据与关联信息\\本地名录R3、带标签近邻};

  \node[model block] (model) at (70,-27) {标注模型\\版本$M_5$};
  \node[output block] (output) at (101,-27) {候选记录\\答案、分数\\证据与版本};

  \draw[convention] (requirement.east) -- node[edge text,above] {约定} ([yshift=6mm]model.west);
  \draw[convention] (demonstration.east) -- node[edge text,above] {示范} ([yshift=2mm]model.west);
  \draw[evidence] (sample.east) -- node[edge text,below] {当前观测} ([yshift=-2mm]model.west);
  \draw[evidence] (support.east) -- node[edge text,below] {外部证据} ([yshift=-6mm]model.west);
  \draw[-{Stealth[length=2mm]},draw=FigureModel,line width=1.0pt] (model.east) -- (output.west);
\end{tikzpicture}

  \caption{一次模型标注中的四类输入。虚线表示对作答方式的约定，实线表示支持当前实例判断的观测或证据；实例标识把照片、地点与检索材料绑定到同一对象。四类输入描述功能组成，不表示所有模型都必须使用同一种提示接口。}
  \label{fig:knowledge-acquisition-input-composition}
\end{figure}

输入组织不是把所有可得信息尽量塞给模型。每增加一段上下文，都要检查它与待处理对象是否对应、是否占用有限输入空间、是否与其他内容矛盾，以及是否暴露了评价时本不应获得的答案。检索过程负责找到候选材料，本节负责将找到的材料与当前任务对应起来；两者混为一谈时，检索错误和组织错误将无法分别诊断。

\subsection{任务要求与输出约定}
\label{sec:knowledge-acquisition-task-output-contract}

任务要求先界定答案的语义。花卉识别要说明使用哪套物种名录、杂交种和别名怎样处理、多张照片是否允许对应同一对象，以及何时应保留“未知”或“不适用”。输出约定再规定样本标识、候选内容、来源动作、原因和证据引用等字段；有可解释的模型得分时另行保存。前者回答“判断什么”，后者回答“怎样交付”。S042使用下列任务说明：

\begin{quote}
  请按规范Q1，仅根据本请求提供的观测和证据，判断S042所指植株的物种。任务范围限于类别表C1中的山茶、茶梅、月季。若作答，只返回其中一个规范名称，不自行补入别名、杂交种或其他类别。

  只有证据足以支持唯一类别时才作答，此时内容为物种名称，语义状态为“已给出”，动作为“作答”。观测不足、材料冲突或只能缩小候选范围时，内容留空，语义状态为“未知”，动作为“弃权”。只有已确认对象不属于本任务范围时，才记为“不适用”，内容和动作均留空。不得臆测不可见部位，也不得把某地出现过的物种直接当作该地每株植物的物种。

  只返回一个JSON对象，字段固定为object、content、semantic、action、reason、evidence，依次表示对象标识、内容、语义状态、来源动作、原因和证据编号列表。空值写为null，reason简述判断依据，不虚构概率分数。后附S017是完整示范，不是本次待作答对象，也不是S042的证据。
\end{quote}

执行状态由调用程序记录，不能让模型通过自报“成功”证明调用已经完成；验证状态也由后续检查设置。这样，模型交付的任务内容与运行日志、核验结论各有归属。

对固定分类接口，同样的规范要落实为配置。例如C1规定索引$1\mapsto$山茶、$2\mapsto$茶梅、$3\mapsto$月季，则向量$(0.55,0.40,0.05)$的最大项对应山茶。若教学接收门槛另设为$0.70$，本次最高分不足，包装程序返回未知并弃权，而不是把未知塞进某个物种索引。固定接口仍只接收它支持的照片或特征，不能直接执行上述自然语言请求；示范和证据的功能须由支持这些信息的外围组件实现。字段重命名可能允许迁移旧记录，类别拆分或语义边界改变则可能要求重新预测或复核。

边界案例比长篇抽象说明更能检验规范。对山茶与茶梅，可以各给典型正例、易混反例和未知例，说明叶缘、花期等证据在什么条件下才参与判断；对实体抽取，则要固定实体类型、原文跨度与规范化名称之间的关系。封闭类别要求模型只在给定集合中选择，开放文本允许产生集合外答案，多解任务还要约定候选之间是否互斥。三种输出空间不能共用一个未经说明的“答案”字段。

\subsection{待处理样本与必要上下文}
\label{sec:knowledge-acquisition-sample-context}

待处理样本是当前要获得候选知识的对象，上下文则是解释该对象所需且在实际使用时可获得的观测。对照片，应保存原图或稳定引用、待判断区域、拍摄时间和对象标识；对文本，应保存原文位置和必要段落；对语音和时序记录，还要标明时间区间。先确定对象，再裁取上下文，可以避免模型回答了相邻对象的问题却仍把结果写回当前样本。

在S042请求中加入下面的当前观测。O042是这一观测片段的编号，不是物种标签：

\begin{quote}
  当前对象：S042，对应植株P042。观测O042：2026年4月12日在教学园地块P03拍摄的花朵局部，仅见粉红色花瓣；叶片、叶缘、完整花形与株形均不可见。地点和日期来自该照片的采集记录。本请求没有其他照片，也没有这株植物的已知物种标签。
\end{quote}

因此，后续说明可以使用地点和日期，却不能声称已经观察到叶缘。补入同一植株的叶片图会改变可用观测；把提示改成“仔细看叶缘”则不会补出缺失信息。

局部信息与整体信息的取舍取决于任务。只截取花瓣能够放大形态细节，却可能丢失叶片和株形；保留整张照片补回全局线索，也可能引入花盆、背景标牌等干扰。长文本切分时，重叠窗口有助于保留跨边界信息，但同一实体可能被重复预测，需要在还原阶段按原文位置合并。时序任务只能使用判断时点之前允许获得的历史，不能把未来观测混入上下文后仍称为在线预测。

多模态输入尤其要核对对象、位置和时间。若照片属于植株A，地点描述却来自附近植株B，模型可能给出语言上合理、对象上错误的答案。缺失模态也不应被静默替换成默认值；输出记录需要说明模型实际看见了哪些信息。最终评价样本的上下文整理规则必须在看见答案前冻结，否则一次“补充必要背景”可能实质上把参照答案泄漏给模型。

% 为首段及GPT-3完整作者边注共同预留空间。
\Needspace{125mm}
\subsection{示范案例与参照答案}
\label{sec:knowledge-acquisition-demonstrations}

\term{示范案例}由已经完成的输入和参照答案组成，它让模型看到任务规范怎样落实到具体实例。对于可通过\term{上下文学习}（in-context learning, ICL）接受示范的语言模型，这种使用方式不要求更新参数；GPT-3工作中的少样本设置就是把任务和示范放入文本交互中完成预测的代表实例\scite{brown2020}
这并不意味着示范对任意模型都有效，也不意味着参数保持不变时输出就稳定。

示范选择要覆盖常见情况、易混边界和弃权情形，并保持字段与正式输入一致。若示范只有清晰正面照片，就没有展示证据不足时如何作答，模型可能仍然给出物种名称。加入未知示范可以说明弃权条件，是否有效仍需试标。S042请求加入下面这条完整教学示范，示范答案按Q1的弃权条件构造：

\begin{quote}
  示范输入：对象S017，植株P017，类别表C1。观测O017：2026年4月10日在教学园地块P08拍摄，仅见失焦的粉红色花瓣局部，叶片和株形不可见；没有辅助证据。按同一规范Q1判断物种。示范输出如下：
\end{quote}
\begin{lstlisting}[style=book code,numbers=none]
{
  "object": "S017", "content": null,
  "semantic": "未知", "action": "弃权",
  "reason": "花瓣局部失焦，缺少区分物种的观测",
  "evidence": ["O017"]
}
\end{lstlisting}

这条示范展示的是如何报告证据不足，不是在告诉模型“S042也是S017”或“两株属于同一物种”。示范理由也须核对，否则规范的错误解释会随示范一起进入输入。

数量、顺序和相似性都是需要试验的设计变量。少样本提示的顺序敏感性研究表明，同一组示范的不同排列可以造成明显性能波动，而且一种模型上的好顺序未必迁移到另一种模型\scite{knowledge-acquisition-prompt-order2022}
因此，应在开发样本上比较候选示范方案并记录版本，不把最终评价样本用于反复挑选示范。示范过多还会挤占当前样本和证据的输入空间，增加调用成本；示范收益必须与这些代价一起评价。

示范与近邻证据的界线在于主要作用。固定的一组案例若用于说明输出规则，它们属于示范；根据当前查询动态检索的相似带标签照片若直接参与答案计算，则属于辅助证据，并在第~\ref{sec:knowledge-acquisition-neighbor-prediction}节形成近邻预测。两者可以使用同一底层样本库，但选择逻辑和效果解释不能混用。

\subsection{辅助证据与关联信息}
\label{sec:knowledge-acquisition-auxiliary-evidence}

模型参数未必保存了当前判断所需的全部事实，外部文档、数据库、参考样本和对象关系可以作为\term{辅助证据}。组织过程需要明确找什么、怎样确认材料对应当前对象、向模型提供哪一段内容，以及结果记录怎样回指原始来源。检索命中只说明某项材料与查询相似或字段匹配，不自动说明材料支持某个答案，更不说明答案已经正确。

以花卉识别为例，可以按照片表示检索带标签参考照片，也可以按地点和花期检索地区植物名录。每条材料要保留来源、时间范围、对象标识和已有标签状态；同名植物、过时名录或重复转载都可能造成错误关联。文本任务中，同一个实体名称也可能指向不同对象，必须借助上下文消歧，而不能把字符串相同当作事实相同。

对S042，检索条件固定为地块P03、日期2026年4月12日，只取覆盖该地点和日期的一条名录片段。把下列内容接在示范之后，请求便已包含四类输入：

\begin{quote}
  辅助证据E03：教学园维护名录R3，第8条，记录时间为2026年4月1日。适用对象是地块P03，时间范围是2026年3月1日至6月30日。原文：“本期地块记录中有山茶和茶梅；此表不是完整物种清单，未建立逐株标识与物种的对应。”该片段已与R3核对，但没有核验P042的物种。请在这一范围内使用E03，只为当前对象S042返回一次结果。
\end{quote}

P042的地点和日期落在E03范围内，但地块记录不能确定一株植物的物种，也不能排除清单外类别。下面是一条符合请求的构造输出：

\begin{lstlisting}[style=book code,numbers=none]
{
  "object": "S042", "content": null,
  "semantic": "未知", "action": "弃权",
  "reason": "O042缺少鉴别特征，E03没有逐株对应",
  "evidence": ["O042", "E03"]
}
\end{lstlisting}

若这是一轮实际完成的调用，程序会把执行状态记为“成功”，来源记为模型$M_5$，取得条件保存规范Q1、类别表C1、O042、S017示范与E03的引用，验证状态先记为“候选”。复核者可以确认弃权符合Q1，此时核验的是动作是否合理，物种内容仍然未知。相反，下面这条输出虽然JSON格式有效，却应在内容检查时拒收：

\begin{lstlisting}[style=book code,numbers=none]
{
  "object": "S042", "content": "山茶",
  "semantic": "已给出", "action": "作答",
  "reason": "R3已经证明P042是山茶",
  "evidence": ["E03"]
}
\end{lstlisting}

拒收原因是它把地块级名录改写成了逐株证明。原始输出仍须保留，内容核验状态记为“已拒绝”；不能因答案被拒绝就把已完成调用改记为执行失败。前一个输出可以按弃权记录接收，后一个不能作为物种标签使用，这一区别来自证据范围，而不是输出字段是否齐全。

检索增强生成（retrieval-augmented generation, RAG）给出了“参数模型加外部非参数记忆”的一种具体实现：检索器从索引中取得文档，生成模型再在检索结果条件下产生文本\scite{knowledge-acquisition-rag2020}
本章只借它说明证据检索和模型作答可以组合，不把某一RAG架构当作所有辅助证据流程的定义。材料如何被找到、是否完整，以及模型是否忠实使用材料，仍是三个不同问题。

扩大检索范围可能提高召回，也会增加噪声、重复信息和调用成本。输入组织阶段应限制候选数量，标明证据与样本的对应关系，并在材料冲突时保留冲突而不是任意覆盖。此处只准备可追溯证据；近邻和图结构怎样用这些关系计算答案，分别由第~\ref{sec:knowledge-acquisition-neighbor-prediction}节和第~\ref{sec:knowledge-acquisition-graph-prediction}节说明，多来源结论怎样融合则由第~\ref{chap:knowledge-fusion}章处理。

\section{候选知识的预测}
\label{sec:knowledge-acquisition-candidate-prediction}

输入准备完成以后，还要辨认候选答案究竟怎样计算。图~\ref{fig:knowledge-acquisition-prediction-mechanisms}用同一个待识别样本比较三种机制：已学映射从模型参数中读取输入输出关系，近邻预测在调用时读取参考样本的答案，图结构预测则让多个已知和未知节点通过连接共同确定结果。三者可以组合，例如用神经模型产生表示，再在表示上检索近邻或构图；分类依据是当前答案的主要计算路径，不是底层是否出现神经网络。

\begin{figure*}[htbp]
  \centering
  % 图4-3：三种预测机制对照，按169 mm全版心宽度设计。
\begin{tikzpicture}[
  x=1mm,
  y=1mm,
  font=\figurefont\footnotesize,
  text=BookInk,
  unknown/.style={circle,draw=FigureInput,fill=white,line width=0.9pt,minimum size=8mm,inner sep=0pt},
  known/.style={rectangle,draw=FigureInput,fill=FigureInput!10,line width=0.9pt,minimum size=8mm,inner sep=0pt},
  compute/.style={book node,draw=FigureModel,fill=FigureModel!8,line width=0.8pt,text width=16mm,minimum height=11mm,inner sep=1mm},
  answer/.style={book node,draw=FigureOutput,fill=FigureOutput!6,line width=0.8pt,text width=13mm,minimum height=11mm,inner sep=1mm},
  relation/.style={draw=BookMuted,line width=0.7pt},
  calculation/.style={-{Stealth[length=1.8mm]},draw=FigureModel,line width=0.9pt},
  annotation/.style={font=\figurefont\scriptsize,text=BookMuted,align=center}
]
  \draw[BookRule,line width=0.7pt] (55,2) -- (55,-52);
  \draw[BookRule,line width=0.7pt] (110,2) -- (110,-52);

  \node[anchor=north west,font=\figurefont\bfseries] at (1,1) {A\quad 已学映射};
  \node[unknown] (direct-query) at (9,-25) {S};
  \node[annotation] at (9,-34) {查询S042};
  \node[compute] (direct-model) at (29,-25) {模型参数\\$f_{\bm{\theta}}$};
  \node[answer] (direct-answer) at (48,-25) {候选\\山茶};
  \draw[calculation] (direct-query.east) -- (direct-model.west);
  \draw[calculation] (direct-model.east) -- (direct-answer.west);
  \node[annotation] at (29,-44) {预测时主要读取参数};

  \node[anchor=north west,font=\figurefont\bfseries] at (57,1) {B\quad 近邻参考};
  \node[unknown] (neighbor-query) at (64,-27) {S};
  \node[annotation] at (64,-37) {查询S042};
  \node[known] (neighbor-a1) at (79,-13) {A};
  \node[known] (neighbor-a2) at (79,-27) {A};
  \node[known] (neighbor-b) at (79,-41) {B};
  \node[compute,text width=13mm] (vote) at (98,-27) {加权\\投票};
  \node[answer,text width=13mm] (neighbor-answer) at (98,-46) {候选\\山茶};
  \draw[relation] (neighbor-query) -- (neighbor-a1);
  \draw[relation] (neighbor-query) -- (neighbor-a2);
  \draw[relation] (neighbor-query) -- (neighbor-b);
  \draw[calculation] (neighbor-a1.east) -- (vote.west);
  \draw[calculation] (neighbor-a2.east) -- (vote.west);
  \draw[calculation] (neighbor-b.east) -- (vote.west);
  \draw[calculation] (vote.south) -- (neighbor-answer.north);
  \node[annotation] at (79,-51) {$A=$山茶，$B=$茶梅};

  \node[anchor=north west,font=\figurefont\bfseries] at (112,1) {C\quad 图结构};
  \node[known] (graph-a) at (119,-14) {A};
  \node[known] (graph-b) at (119,-42) {B};
  \node[unknown] (graph-u1) at (132,-27) {S};
  \node[unknown] (graph-u2) at (144,-15) {?};
  \node[unknown] (graph-u3) at (144,-40) {?};
  \node[answer,text width=13mm] (graph-answer) at (161,-27) {S042\\山茶};
  \draw[relation] (graph-a) -- (graph-u1);
  \draw[relation] (graph-b) -- (graph-u1);
  \draw[relation] (graph-a) -- (graph-u2);
  \draw[relation] (graph-u1) -- (graph-u2);
  \draw[relation] (graph-u1) -- (graph-u3);
  \draw[relation] (graph-b) -- (graph-u3);
  \draw[relation] (graph-u2) -- (graph-u3);
  \draw[calculation] (graph-u1.east) -- (graph-answer.west);
  \node[annotation] at (136,-51) {方形为标签锚点；边无方向};
\end{tikzpicture}

  \caption{三种预测机制怎样使用已有知识。圆形表示答案未知的查询或未标注节点，方形表示带已知答案的参考或锚点；无箭头边表示相似关系，带箭头路径表示计算。图中类别和得分只演示计算机制，未应用Q1的接收约定，不替代前节对S042的弃权判断；二维位置、距离和边权也均为教学示意。}
  \label{fig:knowledge-acquisition-prediction-mechanisms}
\end{figure*}

表~\ref{tab:knowledge-acquisition-prediction-mechanisms}从预测时保留的信息、核心计算和失效条件比较三种机制。它们不是互斥模型族：直接模型可以负责表示，近邻可以提供证据，图推断可以在一批样本间协调答案。组合以后仍需分别检查参数映射、检索关系和图结构，不能只用最终输出掩盖中间依赖。

\begin{table*}[htbp]
  \centering
  \small
  \setlength{\RaggedRightParindent}{0pt}
  \renewcommand{\arraystretch}{1.15}
  \begin{tabularx}{\linewidth}{@{}>{\RaggedRight\arraybackslash}p{20mm}*{4}{>{\RaggedRight\arraybackslash}X}@{}}
    \toprule
    预测机制 & 读取的信息 & 核心计算 & 新样本处理 & 主要依赖与风险 \\
    \midrule
    已学映射 & 参数、接口与当前输入 & 前向计算或解码 & 按接口推断 & 训练覆盖、接口与分布变化 \\
    近邻参考 & 参考答案、表示与索引 & 距离排序与加权 & 检索并聚合 & 距离含义、标签质量与检索遗漏 \\
    图结构 & 边权、标签锚点与节点 & 联立求解或传播 & 接入图并更新解 & 错边、错标签与无锚点分量 \\
    \bottomrule
  \end{tabularx}
  \caption{三种预测机制的使用条件。比较对象是得到当前答案时主要读取的信息和执行的计算，不是模型结构的永久类别；学习表示、检索证据和图推断可以在同一方案中组合。}
  \label{tab:knowledge-acquisition-prediction-mechanisms}
\end{table*}

\subsection{基于已学映射的直接预测}
\label{sec:knowledge-acquisition-direct-prediction}

\term{直接预测}把组织后的输入交给训练完成的模型，主要利用参数中已经学到的输入输出关系得到候选答案。设$\mathsf{u}_i$表示样本$i$连同任务要求、必要上下文和可选辅助内容组成的模型输入。对封闭的类别集合$Y=\{1,\ldots,K\}$，模型可以给出类别得分$s_{\bm{\theta},c}(\mathsf{u}_i)$，再按
\begin{equation}
  \widehat{y}_i
  \in
  \argmax_{c\in Y}s_{\bm{\theta},c}(\mathsf{u}_i)
  \label{eq:knowledge-acquisition-direct-classification}
\end{equation}
形成候选类别。若需要弃权，还要另外规定最大得分、类别间隔或其他条件；得分只有在模型和评价支持时才能解释成概率，$\argmax$结果也不携带答案正确的保证。

直接模型的训练目标也可能来自一组样本，而不是每个样本各有一个标签。第~\ref{sec:rule-based-knowledge-acquisition-relation-matching}节取得的远程监督以同一实体对的提及组为单位，不能给组内每句都贴上正标签。沿用其中两句话：句乙“林舟在论坛上批评了海岬科技的发布流程”仅共同提及两者；句甲“海岬科技任命林舟为首席架构师”表达任职关系。这里的语义对照供读者理解，训练时不提供这两个逐句答案，只提供“组内至少一句表达任职关系”的组级监督。

令$B$为组内句子的索引集合，$Z_j\in\{0,1\}$表示第$j$句是否表达目标关系，模型根据该句及实体标识组成的输入$\mathsf{u}_j$计算$p_j=p_{\bm{\theta}}(Z_j=1\mid\mathsf{u}_j)$。一种简单聚合是\term{噪声或}（noisy-or）：用“一句也没有表达”的补事件计算“至少一句表达”的概率。它假定给定整组输入后，各句的$Z_j$条件独立，且各自的概率只由相应输入$\mathsf{u}_j$决定。组标签记为$y_B\in\{0,1\}$，采用自然对数时，组级预测与二元交叉熵为
\begin{equation}
  \begin{aligned}
    \widehat p_B
    &=1-\prod_{j\in B}(1-p_j),\\
    \ell_B
    &=-y_B\log\widehat p_B
      -(1-y_B)\log(1-\widehat p_B).
  \end{aligned}
  \label{eq:knowledge-acquisition-model-bag-prediction}
\end{equation}
对一个人为构造的正组，按句乙、句甲的顺序设$p_1=0.1$、$p_2=0.8$、$y_B=1$，就有
\[
  \widehat p_B=1-(1-0.1)(1-0.8)=0.82,
  \qquad
  \ell_B=-\log0.82\approx0.19845.
\]
训练时，梯度经聚合式传到句级模型；它拟合的是“至少一句表达”的事件，而不是把两句分别当作确定正例。预测时可以返回组级分数$0.82$，并将句甲作为分数较高的证据候选；该组分数不是句甲或句乙各自的概率，组级标签也不足以唯一确定支撑句。

上述条件独立是一项建模假设。重复转载或高度相似的句子会产生依赖，直接相乘可能高估组级支持。更根本的风险是知识库事实未必在当前提及组中被任何一句表达：若$y_B=1$只是远程监督给出的候选，便须另行接收、降权或复核，不能把它当作已核验的组级事实。上式展示了接受一条组级监督后的训练计算，没有替代这一步质量判断。

回归把输出解释为数值，结构化预测还要在多个局部结果之间满足序列、树或集合约束。开放文本生成则把答案分解为词元序列：
\begin{equation}
  p_{\bm{\theta}}(y_{1:m}\mid\mathsf{u}_i)
  =
  \prod_{r=1}^{m}
  p_{\bm{\theta}}(y_r\mid y_{1:r-1},\mathsf{u}_i).
  \label{eq:knowledge-acquisition-direct-generation}
\end{equation}
解码策略、候选数量和最大长度会改变实际输出；生成物种名称以后，还要映射到类别表并处理别名。能够产生一段通顺解释，只说明模型完成了文本生成，不说明解释中的证据真实，也不说明物种答案可靠。

实际批量处理还包括切分、组批、异常返回和结果还原。长文本的多个窗口可能给同一实体产生冲突答案，图像接口可能因文件损坏而返回空结果，结构化输出也可能语法有效但字段语义错误。原始输出和解析后的候选都应保留，使后续能够区分模型生成错误、解析错误和写回对象错误。增加采样次数或候选数量可以扩大结果集合，也会增加融合和复核负担；多份候选如何综合由第~\ref{chap:knowledge-fusion}章处理。

直接预测也可以共同生成任务、输入和答案，而不只是给现有样本补答案。\term{Self-Instruct}让语言模型基于已有种子任务生成新指令，并为指令构造输入和输出，再过滤无效或相似实例，用保留材料继续微调模型\scite{wang2023selfinstruct}
这里变化的不只是标签数量，训练任务的覆盖和分布也被模型改写。因此，检查流程必须同时覆盖任务是否有效、输入与指令是否相容、答案是否可用以及材料是否过度重复，不能沿用只核对固定样本标签的办法。生成数量增加只表明候选集扩大，不证明任务覆盖或训练效用提高。

\subsection{基于近邻参考的预测}
\label{sec:knowledge-acquisition-neighbor-prediction}

\term{近邻参考预测}在调用时保留一组带已知答案的参考样本，并让与查询相近的参考项参与当前判断。设参考集为$\mathcal{R}=\{(\bm{x}_j,y_j)\}_{j=1}^{n}$，$\phi$把样本映射到比较空间，$d$是该空间上的距离，$N_k(\bm{x})$是查询$\bm{x}$的$k$个近邻索引。对分类任务，一种加权投票为
\begin{equation}
  \begin{aligned}
    \widehat{y}(\bm{x})
    &\in
    \argmax_{c\in Y}
    \sum_{j\in N_k(\bm{x})}
    \alpha_j(\bm{x})\mathbf{1}\{y_j=c\}, \\
    \alpha_j(\bm{x})
    &=
    \frac{w\!\left(d(\phi(\bm{x}),\phi(\bm{x}_j))\right)}
    {\sum_{r\in N_k(\bm{x})}w\!\left(d(\phi(\bm{x}),\phi(\bm{x}_r))\right)}.
  \end{aligned}
  \label{eq:knowledge-acquisition-neighbor-prediction}
\end{equation}
其中$w$是非负权重函数，分母要求至少一个邻居具有正权重。回归可以把指示函数换成参考数值$y_j$，得到加权平均。距离并列、类别平票、无有效邻居和查询与参考完全重复时，都要预先规定处理方式。

经典近邻分类说明了局部参考怎样直接形成预测\scite{coverhart1967}
完整算法与统计性质可回看第~\ref{sec:classification-knn}节。本章关注的是知识依赖：即使$\phi$由神经模型学习，式~\eqref{eq:knowledge-acquisition-neighbor-prediction}仍在预测时读取参考样本的答案，因此不同于只读取参数的直接映射。删除或更正一个参考标签可能立刻改变附近查询的结果，而不需要重新训练表示模型。

花卉照片中的“相近”必须与任务含义一致。若表示主要编码背景和拍摄设备，一张公园中的山茶可能检索到同一地点的月季；若表示捕捉花瓣、叶缘和株形，邻居才更可能提供物种证据。邻居一致也不等于答案可靠：参考集可能来自同一错误来源，类别失衡也会让多数类在局部占据过多位置。稀疏区域没有足够近的参考时，应降低支持程度或弃权，而不是强行把最不远的样本称为“相似”。

参考集扩大以后，存储、索引构建和查询成本都会变化。近似近邻索引可以减少搜索成本，却可能漏掉真正最近项；新增或撤回参考标签还可能要求增量维护索引和缓存。因而每条候选答案最好保留所用参考样本标识、距离、权重和索引版本，后续才能重放式~\eqref{eq:knowledge-acquisition-neighbor-prediction}，并判断变化来自参考知识还是检索实现。

\subsection{基于图结构的预测}
\label{sec:knowledge-acquisition-graph-prediction}

近邻预测通常逐个查询并聚合局部参考，\term{图结构预测}则把已知与未知样本同时放入图中，使多个对象之间的连接共同约束答案。设无向加权图为$\mathcal{G}=(V,E)$，边权$w_{i,j}=w_{j,i}\geq0$表示任务相关的相似关系；节点集合分为已有标签的$L$和待求答案的$U$。对$K$类任务，用$\bm{F}_i\in\mathbb{R}^{K}$表示节点$i$的类别得分，已知节点固定为独热向量$\bm{Y}_i$。一种\term{标签传播}（label propagation）目标是
\begin{equation}
  \min_{\bm{F}:\,\bm{F}_L=\bm{Y}_L}
  \frac{1}{2}\sum_{i,j\in V}w_{i,j}\lVert\bm{F}_i-\bm{F}_j\rVert_2^2
  =
  \min_{\bm{F}:\,\bm{F}_L=\bm{Y}_L}
  \operatorname{tr}\!\left(\bm{F}^{\top}\bm{L}_{\mathcal{G}}\bm{F}\right),
  \label{eq:knowledge-acquisition-graph-energy}
\end{equation}
其中$\bm{L}_{\mathcal{G}}=\bm{D}-\bm{W}$为图Laplacian，$\bm{D}_{i,i}=\sum_jw_{i,j}$。目标让高权重相连节点的得分接近，同时把已有标签作为边界条件固定下来。无向边表示相似关系，不表示一个样本导致另一个样本具有某个类别。

按照$L,U$分块，对未知部分求驻点得到
\begin{equation}
  \bm{L}_{U,U}\bm{F}_U
  =
  -\bm{L}_{U,L}\bm{Y}_L,
  \qquad
  \bm{F}_U
  =
  -\bm{L}_{U,U}^{-1}\bm{L}_{U,L}\bm{Y}_L.
  \label{eq:knowledge-acquisition-graph-solution}
\end{equation}
这个逆矩阵形式只有在$\bm{L}_{U,U}$可逆时成立。对非负无向边权，如果每个含未知节点的连通分量都能沿某条路径到达至少一个已标注节点，固定边界后的能量有唯一解；若某个未知连通分量完全没有标签锚点，它的绝对类别无法由该目标确定，$\bm{L}_{U,U}$也会出现相应奇异性。Gaussian fields与调和函数方法给出了这一构造、矩阵解和随机游走解释\scite{zhu2003}

式~\eqref{eq:knowledge-acquisition-graph-solution}也可以用局部迭代理解。对未知节点$i$，调和条件要求
\[
  \bm{F}_i
  \leftarrow
  \frac{\sum_j w_{i,j}\bm{F}_j}{\sum_j w_{i,j}},
\]
并在每轮把已知节点重新固定为$\bm{Y}_i$。直接近邻只汇总查询附近的已知答案，图迭代则允许信息通过多个未知节点继续传递。传播范围扩大以后，错误种子或跨类别错误连边也会影响更远节点；算法数值收敛不代表图假设正确。

\begin{example}[直接近邻、多步传播与错误连边]
  \label{ex:knowledge-acquisition-graph-propagation}
  只看类别A的一维得分。令A、B两个锚点的得分分别为$f_A=1$和$f_B=0$，查询节点$u$以权重2和1连接A、B，并以权重1连接未知节点$v$；$v$再以权重3连接A。若忽略未知邻居，只按已标直接近邻归一化，得到
  \[
    f_u^{\mathrm{direct}}=\frac{2f_A+f_B}{2+1}=\frac{2}{3}.
  \]
  使用完整图的调和条件时，
  \[
    4f_u=2+f_v,\qquad 4f_v=3+f_u,
  \]
  所以$f_u=11/15$、$f_v=14/15$。多步传播通过$v$把额外的A类锚点信息带到$u$。若再错误地给$v$与B添加一条权重3的跨类边，第二式变为$7f_v=3+f_u$，重算得到$f_u=17/27$、$f_v=14/27$。若接收阈值为$0.7$，完整正确图会接收$u$的A类候选，加入错误边后则不再接收。图~\ref{fig:knowledge-acquisition-prediction-mechanisms}(c)复用了这组节点、边权和结果。
\end{example}

花卉图可以把照片作为节点，把形态表示、拍摄地点或同一植株关系转成边权。加入一条把背景相似但物种不同的照片强连接起来，错误就可能沿图扩散；新增一批照片也会改变邻接关系，使旧节点得分发生变化。构图、稀疏存储和线性方程求解都是成本来源。节点表示和边权可以学习，但本节的输出仍由已有任务标签锚定；第~\ref{chap:data-derived-knowledge-acquisition}章讨论的数据分组只发现尚未赋予物种语义的群组，不能因为都使用图就把二者混为一类。

\section{获取过程的迭代改进}
\label{sec:knowledge-acquisition-iterative-improvement}

一轮获取结束以后，复核结果可能指出不同问题：模型没有学会类别边界，输入缺少叶片证据，检索材料关联错误，或者大量预算花在已经稳定的简单样本上。这些现象不能全部用“再训练一次”处理。迭代开始前应冻结上一轮的数据、模型、输入配置、检索或图版本与输出，随后按成因选择更新模型、更新输入或更新获取安排，并记录哪些旧结果因此失效。

迭代证据也有用途边界。开发验证和过程审计用于选择阈值、输入版本、教师更新方式以及继续、停止或回退；它们一旦参与选择，就不能再充当最终独立验收。图~\ref{fig:knowledge-acquisition-iterative-update}把版本选择与留出验收分开，并在下方面板展开任务教师更新。图中没有从训练损失直接指向“质量改善”，因为目标下降只说明模型更好地拟合当前训练信号。

\subsection{模型的更新}
\label{sec:knowledge-acquisition-model-update}

\term{自训练}（self-training）把模型对未标注样本的部分预测暂时作为训练目标，再用扩充后的训练信号更新模型。设真实标注索引集为$L$，第$t$轮未标注候选集为$U^{(t)}$，模型给出的类别和接收分数分别为$\widehat y_j^{(t)}$与$q_j^{(t)}$。若当前轮采用阈值$\tau^{(t)}$，被接收的伪标签集合可以写成
\begin{equation}
  P^{(t)}
  =
  \left\{
    \left(\bm{x}_j,\widetilde y_j^{(t)},q_j^{(t)}\right)
    :
    j\in U^{(t)},\
    q_j^{(t)}\geq\tau^{(t)},\
    \widetilde y_j^{(t)}=\widehat y_j^{(t)}
  \right\}.
  \label{eq:knowledge-acquisition-pseudo-label-set}
\end{equation}
这里的$\widetilde y_j^{(t)}$是\term{伪标签}：模型提出、经过当前接收规则筛选的临时训练目标，不是真实标签。接收分数也不必然等于正确概率；若其尺度没有校准，阈值只是在该模型版本内部定义筛选规则。

当$P^{(t)}$非空时，一种示意训练目标为
\begin{equation}
  \mathcal{L}^{(t)}(\bm{\theta})
  =
  \frac{1}{|L|}\sum_{i\in L}\ell\!\left(f_{\bm{\theta}}(\bm{x}_i),y_i\right)
  +
  \lambda^{(t)}
  \frac{1}{|P^{(t)}|}
  \sum_{(\bm{x}_j,\widetilde y_j^{(t)},q_j^{(t)})\in P^{(t)}}
  a_j^{(t)}
  \ell\!\left(f_{\bm{\theta}}(\bm{x}_j),\widetilde y_j^{(t)}\right),
  \label{eq:knowledge-acquisition-self-training-objective}
\end{equation}
其中$a_j^{(t)}\geq0$是样本权重，$\lambda^{(t)}\geq0$控制整批伪标签相对于人工标签的影响。若$P^{(t)}$为空，应跳过伪标签项并停止、调整接收规则或补充证据，不能以除零后的任意约定掩盖“本轮没有可接收样本”。真实标签与伪标签分开归一化，使$\lambda^{(t)}$明确控制两类信号的相对量级；若按所有样本统一平均，伪标签数量增加本身就会改变权重。

% 本段末尾两条边注不能跨页续排，需连同段落检查可用高度。
\Needspace{110mm}
基本循环是“预测候选—选择或复核—更新参数—重新预测”。每轮都要保存人工标签、伪标签、样本权重和模型版本，并明确历史伪标签是固定、重估还是允许撤销。一次性加入大量伪标签会迅速扩大覆盖，也可能一次放大系统性错误；分批加入便于观察变化，却增加预测和训练轮次。直接继续训练保留当前优化状态，重新训练更容易隔离历史路径影响，两者都不是无条件更优的安排。早期自训练思想已用于自适应模式识别\scite{scudder1965selftraining}
现代深度模型中的研究也观察到错误伪标签可能被后续训练强化，形成确认偏差\scite{arazo2020confirmation}

任务教师的更新方式决定下一轮目标来自哪个模型。固定教师在多轮中保持不变，便于追溯，却不会吸收学生的新能力；定期替换教师把某个学生版本提升为新教师；随学生更新的教师则在更细时间尺度上变化。定期替换可以使用更大或不同结构的学生，这说明“教师”是角色而不是固定的容量排序，但替换后的教师也可能继承并强化旧错误。

\term{Mean Teacher}采用随学生更新的教师，并以学生参数的指数移动平均（exponential moving average, EMA）更新教师参数\scite{knowledge-acquisition-mean-teacher2017}

第$t$步先用当前学生参数$\bm{\theta}^{(t)}$和教师参数$\overline{\bm{\theta}}^{(t)}$计算损失，再用梯度得到学生参数$\bm{\theta}^{(t+1)}$，最后用更新后的学生参数计算教师参数$\overline{\bm{\theta}}^{(t+1)}$：
\begin{equation}
  \begin{aligned}
    \mathcal{L}_{\mathrm{con}}^{(t)}
      \bigl(\bm{\theta},\overline{\bm{\theta}}^{(t)}\bigr)
    &=
    d\!\left(
      f_{\bm{\theta}}(\widetilde{\bm{x}}_s),
      f_{\overline{\bm{\theta}}^{(t)}}(\widetilde{\bm{x}}_t)
    \right), \\
    \mathcal{J}^{(t)}
      \bigl(\bm{\theta};\overline{\bm{\theta}}^{(t)}\bigr)
    &=
    \mathcal{L}_{\mathrm{sup}}^{(t)}(\bm{\theta})
    +\lambda_{\mathrm{con}}^{(t)}
      \mathcal{L}_{\mathrm{con}}^{(t)}
      \bigl(\bm{\theta},\overline{\bm{\theta}}^{(t)}\bigr), \\
    \bm{\theta}^{(t+1)}
    &=
    \bm{\theta}^{(t)}
    -\eta^{(t)}
      \left.
      \nabla_{\bm{\theta}}\mathcal{J}^{(t)}
      \bigl(\bm{\theta};\overline{\bm{\theta}}^{(t)}\bigr)
      \right|_{\bm{\theta}=\bm{\theta}^{(t)}}, \\
    \overline{\bm{\theta}}^{(t+1)}
    &=
    \alpha\overline{\bm{\theta}}^{(t)}
    +(1-\alpha)\bm{\theta}^{(t+1)}.
  \end{aligned}
  \label{eq:knowledge-acquisition-mean-teacher}
\end{equation}
其中$\mathcal{L}_{\mathrm{sup}}^{(t)}$使用当前批次的真实标签，$\lambda_{\mathrm{con}}^{(t)}\geq0$控制一致性项，$\eta^{(t)}>0$是学生更新步长，$\alpha\in[0,1)$；$\widetilde{\bm{x}}_s$与$\widetilde{\bm{x}}_t$表示可采用不同扰动的同一样本，$d$衡量学生预测与教师目标的一致程度。教师不接收这一步的梯度，更新后的$\overline{\bm{\theta}}^{(t+1)}$在第$t+1$步重新产生目标。原方法的关键是平均模型权重，而不是保存并平均每个样本的历史预测。

例如，在单个标量参数的教学示例中，若$\alpha=0.9$、当前教师参数为$0.6$、更新后学生参数为$1.0$，下一步教师参数为$0.9\times0.6+0.1\times1.0=0.64$。下一轮目标由$f_{0.64}$重新计算，一般不等于$0.9f_{0.6}+0.1f_{1.0}$，因为模型输出对参数通常不是线性的。

上例说明教师参数如何更新。若进一步把这个教师接入式~\eqref{eq:knowledge-acquisition-pseudo-label-set}的阈值伪标签流程，还要重新预测，再决定是否接收候选；硬阈值筛选不是Mean Teacher软一致性训练的必需步骤。为演示这种组合，假设同一候选在更新前由$f_{0.6}$给出接收分数$0.72$，更新后由$f_{0.64}$给出$0.78$，两轮阈值均为$0.75$。该候选在前一轮未进入训练集合，下一轮才以教师预测的类别作为硬伪标签进入$P^{(t+1)}$。这些分数是额外设定的构造值，不能从参数$0.6$和$0.64$直接推得，也不表示模型输出随参数单调变化。

\begin{figure*}[!t]
  \centering
  % 图4-4：候选知识进入下一轮模型更新，按169 mm全版心宽度设计。
\begin{tikzpicture}[
  x=0.96mm,
  y=1mm,
  font=\figurefont\footnotesize,
  text=BookInk,
  loop node/.style={book node,draw=FigureModel,fill=FigureModel!7,line width=0.8pt,text width=20mm,minimum height=12mm,inner sep=1.2mm},
  data node/.style={book node,draw=FigureInput,fill=FigureInput!7,line width=0.8pt,text width=18mm,minimum height=12mm,inner sep=1.2mm},
  audit node/.style={book node,draw=BookAmber,fill=BookAmber!7,line width=0.8pt,text width=22mm,minimum height=12mm,inner sep=1.2mm},
  state node/.style={book node,draw=BookMuted,fill=BookPaper,line width=0.7pt,text width=20mm,minimum height=11mm,inner sep=1.2mm},
  loss node/.style={book node,draw=FigureLoss,fill=FigureLoss!7,line width=0.8pt,text width=15mm,minimum height=11mm,inner sep=1.2mm},
  main flow/.style={knowledge arrow,draw=FigureModel,line width=0.9pt,rounded corners=1mm},
  evidence flow/.style={knowledge arrow,draw=BookAmber,line width=0.8pt,rounded corners=1mm},
  external flow/.style={knowledge arrow,draw=BookMuted,line width=0.8pt,dashed,rounded corners=1mm},
  gradient flow/.style={knowledge arrow,draw=FigureLoss,line width=1.0pt,rounded corners=1mm},
  ema flow/.style={knowledge arrow,draw=BookMuted,line width=0.8pt,dashed},
  edge text/.style={fill=white,inner sep=0.7pt,font=\figurefont\scriptsize,text=BookMuted,align=center}
]
  \draw[BookRule,line width=0.7pt] (111,2) -- (111,-103);

  \node[anchor=north west,font=\figurefont\bfseries] at (1,1) {A\quad 获取与版本决策};
  \node[data node] (data) at (11,-12) {人工标签与\\待处理样本};
  \node[loop node] (model) at (37,-12) {当前模型\\版本$M^{(t)}$};
  \node[loop node] (candidate) at (64,-12) {候选输出\\及原始记录};
  \node[audit node] (selection) at (93,-12) {选择或复核\\接收训练目标};
  \node[state node] (input-change) at (12,-32) {输入修改\\调用安排};
  \node[audit node] (audit) at (63,-31) {开发验证\\过程审计};
  \node[audit node] (decision) at (93,-31) {继续／停止\\／回退？};
  \node[loop node] (update) at (93,-50) {训练更新\\形成$M^{(t+1)}$};
  \node[state node] (previous) at (31,-50) {上一可用版本\\及旧结果};
  \node[state node,draw=FigureOutput] (freeze) at (63,-68) {冻结方案\\模型与输入版本};
  \node[state node,draw=FigureOutput,fill=FigureOutput!6] (holdout) at (63,-84) {留出验收\\独立质量证据};

  \draw[main flow] (data.east) -- (model.west);
  \draw[main flow] (model.east) -- (candidate.west);
  \draw[main flow] (candidate.east) -- (selection.west);
  \draw[evidence flow] (selection.south) -- ++(0,-5) -| (audit.north);
  \draw[evidence flow] (audit.east) -- (decision.west);
  \draw[evidence flow] (decision.south) -- node[edge text,right] {继续} (update.north);
  \draw[main flow] (update.west) -- (78,-50) -- (78,-41) -- (37,-41) -- (model.south);
  \draw[external flow] (input-change.east) -- node[edge text,above] {外部反馈入口} (model.south west);
  \draw[evidence flow] (decision.west) -- (81,-31) |- node[edge text,above,pos=.72] {回退} (previous.north);
  \draw[main flow] (previous.north) -- (31,-41) -- (37,-41);
  \draw[evidence flow] (decision.east) -- (106,-31) |- node[edge text,right,pos=.42] {停止} (freeze.east);
  \draw[main flow] (freeze.south) -- (holdout.north);

  \node[anchor=north west,font=\figurefont\bfseries] at (114,1) {B\quad Mean Teacher};
  \node[loop node,text width=17mm] (student) at (124,-12) {学生参数\\$\bm{\theta}^{(t)}$};
  \node[loop node,text width=17mm] (teacher) at (156,-12) {教师参数\\$\overline{\bm{\theta}}^{(t)}$};
  \node[state node,text width=17mm] (student-output) at (124,-31) {学生预测\\第$t$步};
  \node[state node,text width=17mm] (teacher-output) at (156,-31) {教师目标\\第$t$步};
  \node[data node,text width=10mm] (labels) at (118,-49) {真实\\标签};
  \node[loss node,text width=12mm] (supervised) at (133,-49) {监督\\损失};
  \node[loss node,text width=17mm] (consistency) at (156,-49) {一致性\\损失};
  \node[loop node,text width=18mm] (optimizer) at (137,-66) {优化器更新\\旧参数$+$梯度};
  \node[loop node,text width=17mm] (student-next) at (126,-83) {学生参数\\$\bm{\theta}^{(t+1)}$};
  \node[loop node,text width=17mm] (teacher-next) at (157,-83) {教师参数\\$\overline{\bm{\theta}}^{(t+1)}$};
  \node[state node,text width=17mm,draw=FigureOutput,fill=FigureOutput!6] (next-target) at (157,-100) {教师目标\\第$t+1$步};

  \draw[main flow] (student.south) -- (student-output.north);
  \draw[main flow] (teacher.south) -- (teacher-output.north);
  \draw[main flow] (student-output.south) -- ([xshift=-2mm]supervised.north);
  \draw[main flow] (student-output.east) -- (140,-34) |- (consistency.north west);
  \draw[main flow] (teacher-output.south) -- (consistency.north);
  \draw[main flow] (labels.east) -- (supervised.west);
  \draw[main flow] (student.west) -- (113,-12) |- node[edge text,left,pos=.72] {旧参数} (optimizer.west);
  \draw[gradient flow] (supervised.south) -- node[edge text,left] {梯度} ([xshift=-3mm]optimizer.north);
  \draw[gradient flow] (consistency.south) -- ([xshift=3mm]optimizer.north);
  \draw[main flow] (optimizer.south) -- (student-next.north);
  \draw[ema flow] (student-next.east) -- node[edge text,above] {参数平均} (teacher-next.west);
  \draw[ema flow] (teacher.east) -- (168,-12) |- node[edge text,right,pos=.72] {旧教师} (teacher-next.east);
  \draw[main flow] (teacher-next.south) -- node[edge text,right] {下一轮} (next-target.north);
\end{tikzpicture}

  \caption{候选知识进入下一轮模型更新。面板A中，开发验证和过程审计支持继续、停止或回退；只有继续才触发训练更新，选定方案冻结后才接受留出验收。面板B中，实线表示计算与参数读取，橙色虚线表示仅更新学生的梯度，灰色点线表示教师参数的指数移动平均（EMA）。优化器读取旧学生参数和两项损失的梯度；EMA读取新学生参数与旧教师参数，得到下一步教师后再重新预测。}
  \label{fig:knowledge-acquisition-iterative-update}
\end{figure*}

Mean Teacher的教师输出服务于既定任务的一致性目标，因此在本章属于任务教师；第~\ref{sec:data-derived-knowledge-cross-view-prediction}节还会讨论产生表示目标的目标分支或EMA教师。二者都可能采用参数移动平均，却不能只凭“教师—学生”名称归为同一种知识。无论采用硬伪标签还是软一致性目标，初始错误、简单样本过量和类别覆盖恶化都不会因教师更新自动消失；接收规则、类别覆盖、人工标签表现和资源开销需要同时监测。

定期替换教师的代表是Noisy Student：先由教师为未标注样本生成伪标签，再用带噪学生训练过程结合真实标签与伪标签，最后把学生提升为下一轮教师\scite{knowledge-acquisition-noisy-student2020}

\subsection{输入内容的更新}
\label{sec:knowledge-acquisition-input-update}

复核发现错误以后，先判断模型是否看到了完成任务所需的信息。相似花卉持续混淆，可能是类别边界描述含糊，也可能是照片只包含花瓣、示范没有反例，或者检索结果只按背景相似匹配。四种成因分别对应任务要求、样本上下文、示范案例和辅助证据，修改对象不同，受影响的旧输出也不同。

输入更新应写成可比较的版本差异。例如，规范版本$r_2$在$r_1$基础上补充“叶缘锯齿和花期只能在对应部位或时间信息可见时使用”；示范版本$d_2$加入山茶与茶梅的易混反例；样本版本$x_2$为原照片补入同一植株的叶片图；证据版本$e_2$则替换一条过时名录。每项变化都要记录由哪些错误样本触发、适用于哪些对象，以及哪些旧结果因输入改变必须重算。

为了辨认收益来自哪里，可以保持模型参数不变，只替换一种输入成分进行局部对照。补充区别描述后改善，说明任务边界表达可能是瓶颈；换入反例后改善，说明示范覆盖起作用；加入叶片图后改善，则表明原观测不足。但这些判断仍限于开发样本，反复针对同一小批错误调整会过拟合输入设计。更新还可能挤占上下文空间、引入互相矛盾的说明或增加检索成本，不能只统计修正了多少已知错误。

模型参数更新与输入更新要分别记账。若两者同时改变，即使结果改善，也无法判断作用来自新参数还是新证据；必要时可以采用分阶段试验。规范、示范和证据版本一旦冻结，才进入独立评价。第~\ref{chap:knowledge-quality-assessment-correction}章负责检验版本变化是否真正改善候选知识质量，本节只规定如何把错误诊断转成可追溯的输入修改。

\subsection{获取安排的更新}
\label{sec:knowledge-acquisition-schedule-update}

模型和输入不一定要对所有样本同步刷新。\term{获取安排}规定下一轮处理哪些样本、调用多少次、是否补充证据或转交人工，以及采用全量还是局部重算。安排时应分别查询内容、核验情况、执行阶段和版本依赖，不能把“已确认可用”“需要重算”“证据不足”“等待人工”列成互斥的执行状态。例如S042可以同时保持物种未知，并等待人工补充观测。

执行状态沿用共同协议：尚未发起请求为“未请求”，请求已排队但未开始为“未处理”，已经开始而尚未提交为“处理中”，操作完成为“成功”，因异常终止而未能完成为“失败”。成功弃权仍是成功执行，调用失败则没有主动弃权这一动作。需要重算的是对旧版本依赖的判断，不是新的执行状态；应为新请求单独记录进度，并明确当前使用哪个版本，避免把旧答案与新答案同时视为现行结果。

影响范围由变更对象决定。模型参数整体更新时，所有依赖旧模型直接预测的样本原则上都可能变化；只更正某个类别映射时，可以局部刷新涉及该映射的记录；外部名录更新只影响引用该来源或相应时间范围的候选；近邻参考标签撤回会影响检索到该参考的查询，图中一条关键边变化则可能传到整个连通区域。局部刷新要有可计算的依赖记录，否则节省调用的同时会留下不一致结果。

同样的预算也可以用于不同目的。处理新样本扩大覆盖，重复生成主要暴露随机输出波动，复核困难样本提供错误证据，刷新旧答案处理版本变化，转交人工则补入模型当前缺少的独立信息。重复生成若没有改变模型、输入或证据，只产生高度相关的答案，轮数增加并不等于知识增加。表~\ref{tab:knowledge-acquisition-iteration-actions}将常见反馈与可调整安排对应起来；表中动作是诊断后的选择，不是看到某种现象就必须执行的规则。

\begin{table*}[htbp]
  \centering
  \small
  \setlength{\RaggedRightParindent}{0pt}
  \renewcommand{\arraystretch}{1.15}
  \begin{tabularx}{\linewidth}{@{}>{\RaggedRight\arraybackslash}p{22mm}*{4}{>{\RaggedRight\arraybackslash}X}@{}}
    \toprule
    反馈情形 & 定位范围 & 可选动作 & 额外投入 & 过程检查 \\
    \midrule
    模型更新 & 旧参数或旧表示的依赖项 & 全量或局部刷新 & 训练、重建与推断 & 损失、覆盖与错误类型 \\
    规范变化 & 语义或字段变化的记录 & 重映射、重算或复核 & 规范评审与重算 & 边界案例与新增歧义 \\
    证据缺失 & 缺少必要信息的对象 & 补观测、检索或转人工 & 采集与来源核验 & 对象对应、弃权与错误 \\
    局部错误集中 & 特定类别、来源或图区域 & 定向复核与局部刷新 & 诊断、重建与复核 & 局部收益与其他区域退化 \\
    收益停滞 & 多轮重复的结果 & 停止、回退或换来源 & 比较与回退恢复 & 边际收益、覆盖与成本 \\
    \bottomrule
  \end{tabularx}
  \caption{反馈类型与可调整的获取安排。开发验证和过程审计用于选择继续、停止或回退；参与选择的数据不再充当方案冻结后的最终独立验收证据。}
  \label{tab:knowledge-acquisition-iteration-actions}
\end{table*}

表中的范围要落实到版本依赖：模型更新可能连带改变表示、索引和图，规范变化可能涉及类别映射与解析器，证据更新则需定位原始来源及受影响记录。对局部纠正，应同时检查目标区域和其他区域，防止修正沿参考关系或图边传播后引入新错。

在花卉识别批次中，可以把模型更新后受影响的物种全部重算，为证据缺失照片补拍叶片，对仍无法消除歧义的样本转交人工。若新版本只在已反复调试的样本上改善，而新的开发批次没有收益，应该停止或回退；结果变化越来越小也不等于质量越来越高，它可能只是模型反复产生同一错误。人工查询方法由第~\ref{chap:knowledge-acquisition-human}章展开，跨来源资源分配由第~\ref{chap:knowledge-fusion}章处理，本节只负责模型获取流程中的执行顺序和版本一致性。

\section{本章小结}
\label{sec:knowledge-acquisition-model-based-summary}

基于模型的知识获取不是“选择一个模型并运行”这么简单。完整记录要能够回答四个问题：使用了哪个经过何种准备的模型，模型实际看到了哪些任务要求、观测、示范和证据，当前答案通过已学映射、近邻参考还是图结构怎样计算，以及后续反馈改变了模型、输入还是获取安排。花卉识别中的每个物种候选只有连同这些依赖一起保存，才可能被复核、刷新或撤回。

三种预测机制读取的知识位置不同。直接预测主要依赖模型参数，适合以统一接口批量处理新样本，但受到训练覆盖和输出约束影响；近邻预测在调用时读取参考答案，便于追溯局部依据，却依赖任务相关的表示、距离和参考覆盖；图结构预测利用成批对象之间的连接传播标签，能够跨越直接邻居，也会把错误种子和错误连边传播得更远。三者可以组合，组合不消除各自的假设与维护成本。

任务教师、伪标签、Mean Teacher和Self-Instruct承担不同职责。任务教师是提供任务目标的模型角色，伪标签是被暂时用于训练的一条候选输出，Mean Teacher用学生参数的移动平均产生一致性目标，Self-Instruct则共同生成任务、输入和答案以扩充训练材料。把这些对象分开以后，模型更新、候选生成与材料构造的边界才清楚，也不会把教师名称误当成质量证明。

模型规模更大、上下文更长、参考更多或生成轮次更多，都不能单独证明知识获取更有效。是否继续迭代，要看新证据是否改善已声明范围内的任务结果，错误是否只是重复，以及新增收益能否抵偿训练、检索、推断和复核成本。本章得到的仍是具有任务含义的候选知识；多份结果的综合由第~\ref{chap:knowledge-fusion}章处理，可靠性和冻结方案后的独立证据由第~\ref{chap:knowledge-quality-assessment-correction}章处理，从数据内部构造表示或分组目标则进入第~\ref{chap:data-derived-knowledge-acquisition}章。
